\documentclass[12pt,a4paper]{article}

% Packages for math, code, and formatting
\usepackage{amsmath, amssymb, amsthm}
\usepackage{algorithm}
\usepackage{algpseudocode}
\usepackage{listings}
\usepackage{xcolor}

% Code style for C++
\lstdefinestyle{cpp}{
  language=C++,
  basicstyle=\ttfamily\small,
  keywordstyle=\color{blue},
  commentstyle=\color{green!50!black},
  stringstyle=\color{orange},
  numbers=left,
  numberstyle=\tiny,
  stepnumber=1,
  numbersep=5pt,
  frame=single,
  breaklines=true,
  tabsize=2,
  showstringspaces=false
}

% Theorem environments (for invariants, correctness proofs, etc.)
\newtheorem{theorem}{Teorema}
\newtheorem{lemma}{Lema}
\newtheorem{invariant}{Invariante}
\newtheorem{proposition}{Proposição}

\title{Relatório MC521 - Contest 1}
\author{Sofia Valverde Villas Bôas - 252974}

\begin{document}
\maketitle

\section*{B - P02}

Nesse problema, queremos saber o número máximo de moedas
que é possível conseguir dadas as regras do enunciado relacionada
ao aperto dos botões. 

\subsection*{Solução}

Para a solução, testamos as três possibilidades: apertar os dois botões,
apertar o primeiro botão duas vezes ou apertar o segundo botão duas vezes.
Para cada resultado, calculamos o número de moedas, considerando
que se apertamos o mesmo botão duas vezes, o número de moedas que
ele dá é o número de moedas do botão na primeira vez e o número de moedas
do botão - 1.

Após isso, basta pegar o máximo entre os três resultados e imprimir o resultado.

\section*{H - P08}

Para esse problema, queremos saber a quantidade de pontos feitos por 
Antonia e David ao final do jogo. A cada rodada, ambos cada um rola um dado
e o jogador que tiver a menor pontuação perde pontos equivalentes à pontuação do dado do outro jogador.
Se ambos os jogadores tiverem a mesma pontuação, ninguém perde pontos.

\subsection*{Solução}

Para a solução, lemos o número de rodadas e, para cada rodada, lemos a pontuação de Antonia e David nos dados.

Em seguida, comparamos as pontuações. Se a pontuação de Antonia for maior que a de David, David perde pontos equivalentes à pontuação do dado de Antonia.
Se a pontuação de David for maior que a de Antonia, Antonia perde pontos equivalentes à pontuação do dado de David. 
Caso contrário, ninguém perde pontos.

Por fim, imprimimos a pontuação final de Antonia e David, expressas nas variáveis \texttt{antonia} e \texttt{david}, respectivamente.
\end{document}